% !TeX program = lualatex
% !TeX encoding = utf8
% !TeX spellcheck = uk_UA
% !TeX root = ../ElMagConspect.tex

%=========================================================
\chapter{Енергія магнітного поля}\label{\currfilebase}
\graphicspath{{\currfilebase/Pictures}}
%=========================================================
\epigraph{Магнітне поле не виконує роботи над зарядами — і все ж саме воно зберігає енергію, завдяки якій ця робота стає можливою. Енергія --- не там, де діє сила, а там, де народжується поле.}{Переосмислення ідей сучасної фізики}

На відміну від енергії електростатичного поля, яку можна ввести відразу після закону Кулона (просто обчисливши роботу перенесення зарядів проти кулонівської сили), з магнітною енергією так просто не вийде: магнітна складова сили Лоренца, з якою поле діє на заряд, роботи не виконує (гл.~\ref{MagFieldVacuum}) --- за жодних обставин. Виникає природне запитання: якщо магнітне поле роботи не виконує, звідки ж тоді береться енергія, яку воно, як виявиться, запасає?

Відповідь дає ЕРС \emphx{самоіндукції} --- окремий випадок закону електромагнітної індукції (гл.~\ref{ElMagInduction}): джерело, що підтримує струм у контурі, витрачає частину роботи не на подолання магнітної сили (це неможливо), а на подолання саме цієї ЕРС, яку контур наводить сам у собі через власне змінюване поле. Тому послідовно ввести енергію магнітного поля можна лише \emphx{після} того, як встановлено закон індукції і поняття самоіндукції; сам механізм --- чому саме тут магнітна заборона не діє і хто насправді виконує цю роботу --- ми розберемо детальніше в цьому розділі.

%На відміну від енергії електростатичного поля, яку можна було ввести відразу після встановлення закону Кулона (просто обчисливши роботу
%перенесення зарядів проти кулонівської сили), енергію магнітного поля так само \emphx{прямо} ввести не вдається: саме магнітне поле, за силою
%Лоренца, роботи над зарядами не виконує (гл.~\ref{MagFieldVacuum}), а отже, обчислення роботи струмів проти сили Ампера саме собою нічого не
%скаже про енергію, накопичену в полі. Втім, парадокс тут лише позірний: у гл.~\ref{ElMagInduction} (розд.~\nameref{sec:Work_current}) уже було з'ясовано на мікроскопічному рівні, звідки насправді береться ця робота, попри
%заборону Лоренца, --- заряди в рухомому провіднику мають, крім дрейфової швидкості вздовж нього, ще й швидкість самого провідника, і саме
%складова сили Лоренца, пов'язана з нею, гальмує струм і вимагає компенсації сторонньою ЕРС. Джерело, що підтримує струм у контурі, витрачає
%частину своєї роботи не на подолання магнітної сили як такої, а на подолання саме цієї ЕРС \emphx{самоіндукції} --- тобто ЕРС, яку в контурі
%наводить його ж власне, змінюване з часом магнітне поле. Але ЕРС самоіндукції --- окремий випадок явища електромагнітної індукції. Тому
%послідовно ввести енергію магнітного поля можна лише \emphx{після} того, як встановлено закон електромагнітної індукції і поняття самоіндукції
%(гл.~\ref{ElMagInduction}): без них просто нема з чого обчислювати цю енергію.


%% --------------------------------------------------------
\section{Енергія магнітного поля}
%% --------------------------------------------------------

%% --------------------------------------------------------
\subsection*{Енергія контуру зі струмом}
%% --------------------------------------------------------

Провідник зі струмом створює магнітне поле, яке з'являється й зникає разом зі струмом; \emphx{магнітне поле --- носій енергії}, що дорівнює роботі
джерела струму на його створення. Розглянемо роботу джерела ЕРС $\emf$ при замиканні ключа в колі з котушкою індуктивності $L$ й резистором
$R$, при перенесенні заряду $dq = I\, dt$:
\begin{equation*}
	\delta A = \emf\, dq = I^2R\, dt - \emf_\text{сі}I\, dt = \underbrace{I^2R\, dt}_\text{теплота} +
	\underbrace{\frac1{c^2}LI\, dI}_{\text{енергія поля}}.
\end{equation*}
Частина роботи джерела йде на теплоту (закон Джоуля--Ленца), а решта --- на створення магнітного поля. Інтегруючи другий доданок від нульового
струму до кінцевого значення $I$, отримуємо енергію магнітного поля контуру:
\begin{equation}\label{eq:MagEnergyFormula}
	\tcbhighmath{
		W = \frac1{c^2}\int\limits_0^I LI\, dI = \frac{LI^2}{2c^2} = \frac{I\Phi}{2c} = \frac{\Phi^2}{2L}.
	}
\end{equation}

%% --------------------------------------------------------
\subsection*{Об'ємна густина енергії магнітного поля: соленоїд}
%% --------------------------------------------------------

Знайдемо енергію магнітного поля, використовуючи характеристики самого поля, на прикладі однорідного поля довгого соленоїда з проникністю $\mu$
(рис.~\ref{pic:SolenoidCrossSection}). Індуктивність соленоїда --- за формулою~\eqref{eq:SolenoidInductance}, $L = 4\pi\mu N^2 S/\ell$, а поле
всередині --- $B = 4\pi\mu NI/(c\ell)$ (гл.~\ref{MagFieldMedia}). Підставляючи в~\eqref{eq:MagEnergyFormula}:

\begin{SCfigure}[2][h!]\centering
	\localinput{solenoid-cross-section.tikz}
	\caption{Поздовжній переріз довгого соленоїда: однорідне поле всередині, струм у витках
		обмотки\label{pic:SolenoidCrossSection}}
\end{SCfigure}

\begin{equation*}
	W = \frac1{c^2}\frac{LI^2}2 = \frac1{c^2}\cdot \frac{(4\pi\mu NI)^2 S\ell}{8\pi\mu \ell^2} = \frac{B^2}{8\pi\mu}\, V = \frac{BH}{8\pi}\, V,
\end{equation*}
де $V=S\ell$ --- об'єм соленоїда (об'єм, зайнятий полем), а в останній рівності використано $H=B/\mu$. Звідси об'ємна густина енергії магнітного
поля:
\begin{equation}\label{eq:MagEnergyDensitySolenoid}
	\tcbhighmath{
		w = \frac{BH}{8\pi}.
	}
\end{equation}

\begin{Attention}
	Так само, як і для електричного поля (гл.~\ref{ElEnergy}), \emphx{магнітна енергія зосереджена саме в тому об'ємі простору, де присутнє
		магнітне поле} --- у випадку соленоїда це об'єм усередині нього, а не сам дріт обмотки.
\end{Attention}

%% --------------------------------------------------------
\subsection*{Об'ємна густина енергії магнітного поля: загальний випадок}
%% --------------------------------------------------------

Формулу~\eqref{eq:MagEnergyDensitySolenoid} було отримано лише для окремого випадку однорідного поля соленоїда. Покажемо, спираючись на
загальний закон збереження енергії, що вона (як і аналогічна формула для електричного поля) насправді \emphx{універсальна}.

Саме магнітне поле, нагадаємо, роботи над зарядами не виконує (гл.~\ref{MagFieldVacuum}); всю роботу над носіями струму виконує \emphx{вихрове
	електричне поле} $\Efield$, яке виникає (за законом індукції) щоразу, коли магнітне поле змінюється. За законом Джоуля--Ленца
(гл.~\ref{SteadyCurrent}), робота поля $\Efield$ над струмами $\vect j$ в одиниці об'єму за одиницю часу дорівнює $\vect j\cdot\Efield$, тому за
час $dt$ в об'ємі $V$
\begin{equation*}
	\delta A = dt\iiint\limits_V \vect j\cdot\Efield\, dV.
\end{equation*}
З теореми про циркуляцію для $\Hfield$ у диференціальній формі (гл.~\ref{MagFieldMedia}), $\vect j = \dfrac{c}{4\pi}\rot\Hfield$, тому
\begin{equation*}
	\delta A = \frac{c\, dt}{4\pi}\iiint\limits_V \Efield\cdot\rot\Hfield\, dV.
\end{equation*}
Скориставшись тотожністю $\divg(\Efield\times\Hfield) = \Hfield\cdot\rot\Efield - \Efield\cdot\rot\Hfield$:
\begin{equation*}
	\delta A = \frac{c\,dt}{4\pi}\iiint\limits_V \left[\Hfield\cdot\rot\Efield - \divg(\Efield\times\Hfield)\right]dV.
\end{equation*}
Другий доданок, за теоремою Остроградського--Гаусса, дає потік вектора $\Efield\times\Hfield$ крізь нескінченно віддалену поверхню, який дорівнює
нулю (поля обмеженої системи спадають на нескінченності, гл.~\ref{ElEnergy}):
\begin{equation*}
	\iiint\limits_V \divg(\Efield\times\Hfield)\, dV = \oiint\limits_{S\to\infty} (\Efield\times\Hfield)\cdot d\vect S = 0.
\end{equation*}
У першому ж доданку скористаємось щойно виведеним законом електромагнітної індукції~\eqref{eq:CurlEDiff}, $\rot\Efield =
-\dfrac1c\dfrac{\partial\Bfield}{\partial t}$:
\begin{equation*}
	\delta A = \frac{c\,dt}{4\pi}\iiint\limits_V \Hfield\cdot\rot\Efield\, dV = -\frac{dt}{4\pi}\iiint\limits_V \Hfield\cdot\frac{\partial\Bfield}{\partial
	t}\, dV.
\end{equation*}
Ця робота --- це енергія, яку електричне поле \emphx{забирає} в струмів (щоб віддати її магнітному полю); тому приріст магнітної енергії $dW = -\delta A$:
\begin{equation*}
	dW = \frac1{4\pi}\iiint\limits_V \Hfield\cdot d\Bfield\, dV.
\end{equation*}
Для лінійного середовища ($\Bfield=\mu\Hfield$) маємо $\Hfield\cdot d\Bfield = \tfrac12 d(\Hfield\cdot\Bfield)$, і інтегрування дає остаточний,
\emphx{цілком загальний} результат:
\begin{equation}\label{eq:MagEnergyDensityGeneral}
	\tcbhighmath{
		W = \iiint\limits_V \frac{\Bfield\cdot\Hfield}{8\pi}\, dV, \qquad w = \frac{\Bfield\cdot\Hfield}{8\pi} = \frac{\mu H^2}{8\pi}.
	}
\end{equation}

\begin{Attention}
	Це виведення --- дослівна магнітна паралель до виведення густини енергії електричного поля (гл.~\ref{ElEnergy}), з тією відмінністю, що там
	роль джерела енергії відігравало саме електростатичне поле (яке виконувало роботу над зарядами напряму), а тут --- \emphx{вихрове} електричне
	поле, породжене зміною $\Bfield$ (адже саме магнітне поле роботи виконувати не може). Формула~\eqref{eq:MagEnergyDensityGeneral}, на відміну
	від~\eqref{eq:MagEnergyDensitySolenoid}, справедлива для \emphx{довільного}, як завгодно неоднорідного магнітного поля --- не лише всередині
	соленоїда.
\end{Attention}


%% --------------------------------------------------------
\section{Взаємна індукція}
%% --------------------------------------------------------

%% --------------------------------------------------------
\subsection*{Коефіцієнти взаємної індуктивності}
%% --------------------------------------------------------

Нехай два нерухомі контури 1 і 2 розташовані близько один до одного (рис.~\ref{pic:TwoContours}). Якщо в контурі 1 тече струм $I_1$, він створює
крізь контур 2 магнітний потік, пропорційний $I_1$:
\begin{equation}\label{eq:MutualFlux21}
	\Phi_{21} = \frac1c L_{21}I_1;
\end{equation}
аналогічно, струм $I_2$ у контурі 2 створює крізь контур 1 потік
\begin{equation}\label{eq:MutualFlux12}
	\Phi_{12} = \frac1c L_{12}I_2.
\end{equation}
Коефіцієнти пропорційності $L_{12}$ і $L_{21}$ називають \emphx{коефіцієнтами взаємної індуктивності} контурів; на відміну від власної індуктивності
$L$, це \emphx{алгебраїчні} величини, які можуть мати будь-який знак (залежно від взаємної орієнтації контурів), а також дорівнювати нулю (якщо
контури взаємно перпендикулярні так, що жоден потік іншого через себе не пронизує).

\begin{SCfigure}[1.25][h!]\centering
	\localinput{12contours.tikz}
	\caption{Два нерухомі контури, магнітно пов'язані один з одним\label{pic:TwoContours}}
\end{SCfigure}

З урахуванням взаємної індукції закони Ома для двох контурів набувають вигляду
\begin{equation}\label{eq:MutualOhm}
	\begin{cases}
		I_1R_1 = \emf_1 - \dfrac1{c^2}\dfrac{d}{dt}\left(L_{11}I_1\right) - \dfrac1{c^2}\dfrac{d}{dt}\left(L_{12}I_2\right), \\[1.5ex]
		I_2R_2 = \emf_2 - \dfrac1{c^2}\dfrac{d}{dt}\left(L_{21}I_1\right) - \dfrac1{c^2}\dfrac{d}{dt}\left(L_{22}I_2\right),
	\end{cases}
\end{equation}
де $L_{11}\equiv L_1$, $L_{22}\equiv L_2$ --- власні індуктивності контурів. Той факт, що зміна струму в одному контурі неминуче породжує ЕРС і в
іншому, і називають \emphx{явищем взаємної індукції}.

\begin{Example}{Приклад: дві обмотки на спільному осерді}

Довгий парамагнітний циліндр об'ємом $V$ несе дві обмотки (одна біля одної): перша з густиною намотування $n_1$ витків на одиницю довжини, друга ---
$n_2$; знайдемо їхню взаємну індуктивність, нехтуючи крайовими ефектами.

\begin{center}
    \includegraphics[width=0.65\linewidth]{two_coils}
\end{center}

Оскільки $L_{21}=c\Phi_2/I_1$, потрібно створити струм $I_1$ в обмотці 1 і обчислити повний потік крізь усі $N_2 = n_2\ell$ витків обмотки 2:
$\Phi_2 = N_2 B_1 S$, де $B_1 = \dfrac{4\pi}{c}\mu n_1 I_1$ (гл.~\ref{MagFieldMedia}), тому
\begin{equation}\label{eq:MutualExample}
	L_{21} = 4\pi\mu\, \frac{N_1N_2}{\ell}S,
\end{equation}
і, обчисливши аналогічно $L_{12}$, отримуємо той самий результат: $L_{12}=L_{21}=4\pi\mu N_1N_2S/\ell$.
\end{Example}

\begin{Attention}
	У цьому прикладі $L_{12}=L_{21}=\sqrt{L_1L_2}$. Ця гарна рівність --- наслідок двох спрощень: (1)~ми знехтували розсіюванням магнітного потоку
	поза осердям, і (2)~проникність осердя $\mu$ не залежить від величини поля. Для феромагнітного осердя, де $\mu=\mu(H)$, і за умови $N_1\neq
	N_2$, обмотки при однакових струмах створюють у осерді \emphx{різні} поля, тож вирази для $L_{21}$ і $L_{12}$ міститимуть різні (польозалежні)
	значення $\mu$, і рівність $L_{21}=L_{12}$, узагалі кажучи, порушиться.
\end{Attention}

%% --------------------------------------------------------
\subsection*{Теорема взаємності}
%% --------------------------------------------------------

У щойно розглянутому прикладі коефіцієнти $L_{12}$ і $L_{21}$ обчислювалися двома геометрично різними шляхами, а зійшлися в одному й тому самому
значенні. Це не випадковість прикладу з високосиметричним осердям: зараз буде показано, що рівність $L_{12}=L_{21}$ --- загальний закон, який
виконується для \emphx{будь-якої} пари контурів довільної форми (за єдиної умови --- лінійності середовища, $\mu\neq\mu(H)$, як і вказано вище). Ця рівність, крім самостійного значення, --- ключ до наступного кроку: саме вона дозволить записати енергію довільної системи
зв'язаних контурів компактною формулою.

Означення $L_{21}=c\Phi_2/I_1$ і $L_{12}=c\Phi_1/I_2$ описують, на перший погляд, дві геометрично зовсім різні задачі: обчислення потоку від
контуру 1 крізь контур 2, і навпаки. \emphx{Теорема взаємності} стверджує, що результати цих двох розрахунків завжди --- за
будь-якої форми й взаємного розташування контурів --- точно збігаються:
\begin{equation}\label{eq:MutualEqual}
	\tcbhighmath{
		L_{12} = L_{21}.
	}
\end{equation}
Інакше кажучи: потік, який струм $I$ контуру 1 наводить крізь контур 2, дорівнює потоку, який той самий за величиною струм $I$, потікши тепер
контуром 2, навів би крізь контур 1. Завдяки цій рівності немає потреби розрізняти $L_{12}$ і $L_{21}$ --- можна просто говорити про
\emphx{взаємну індуктивність двох контурів} $L_{12}$. Диференціюючи рівність потоків за часом, отримуємо еквівалентне ЕРС-формулювання: ЕРС,
яку змінний струм $I(t)$ контуру 2 наводить у контурі 1, у кожен момент часу дорівнює ЕРС, яку той самий струм $I(t)$, потікши контуром 1, навів
би в контурі 2,
\begin{equation*}
	-\frac{L_{12}}{c^2}\frac{dI}{dt} = -\frac{L_{21}}{c^2}\frac{dI}{dt}.
\end{equation*}

\subsubsection*{Доведення}

Теорема взаємності аж ніяк не очевидна із самого означення коефіцієнтів --- порівнювати довелося б два геометрично різні розрахунки. Доведемо її
одразу для системи довільної кількості $N$ нерухомих контурів зі струмами $I_1,I_2,\dots,I_N$, скориставшись тим, що магнітна енергія такої
системи --- це функція \emphx{стану}: вона залежить лише від миттєвих значень струмів,
\begin{equation}\label{eq:EnergyAsFunctionOfCurrents}
	W = W(I_1,I_2,\dots,I_N),
\end{equation}
а не від того, у якій послідовності й з якою швидкістю ці струми було встановлено --- так само, як енергія одного контуру в
формулі~\eqref{eq:MagEnergyFormula} не залежить від того, наскільки швидко замикали ключ.

Магнітний потік крізь $i$-й контур --- від власного струму та від струмів усіх інших контурів --- запишемо у вигляді
\begin{equation}\label{eq:FluxThroughIthLoop}
	\Phi_i = \frac1c\sum_{k=1}^{N} L_{ik}I_k,
\end{equation}
де коефіцієнти $L_{ik}$ (сталі, тобто не залежні від струмів) поки що \emphx{не} припускаємо симетричними --- $L_{ik}\equiv\Phi_i/I_k$ при
відсутності струму в решті контурів, і саме симетрію $L_{ik}=L_{ki}$ належить довести. Зміна струмів на $dI_1,\dots,dI_N$ наводить у кожному
контурі ЕРС індукції $\emf_i^\text{інд}=-\frac1c\dfrac{d\Phi_i}{dt}$; щоб підтримати задані струми, джерела виконують проти неї роботу, яка йде
виключно на приріст магнітної енергії системи (як і при виведенні формули~\eqref{eq:MagEnergyFormula} для одного контуру):
\begin{equation}\label{eq:EnergyDifferential1}
	dW = \frac1c\sum_{i=1}^{N} I_i\,d\Phi_i = \frac1{c^2}\sum_{i,k=1}^{N} L_{ik}I_i\,dI_k.
\end{equation}
З іншого боку, оскільки $W$ --- функція незалежних змінних $I_1,\dots,I_N$, її диференціал завжди можна записати як $dW=\sum_k\dfrac{\partial
W}{\partial I_k}\,dI_k$. Порівнюючи це із~\eqref{eq:EnergyDifferential1}, знаходимо частинну похідну енергії:
\begin{equation}\label{eq:PartialDerivatives}
	\frac{\partial W}{\partial I_k} = \frac1{c^2}\sum_{i=1}^{N} L_{ik}I_i.
\end{equation}
Продиференціюємо цей вираз ще раз, за $I_i$:
\begin{equation*}
	\frac{\partial^2 W}{\partial I_i\,\partial I_k} = \frac1{c^2}L_{ik}.
\end{equation*}
Виконуючи диференціювання в оберненому порядку --- спершу за $I_k$, потім за $I_i$, --- за тим самим міркуванням отримаємо
\begin{equation*}
	\frac{\partial^2 W}{\partial I_k\,\partial I_i} = \frac1{c^2}L_{ki}.
\end{equation*}
Але мішана похідна будь-якої гладкої функції не залежить від порядку диференціювання: $\dfrac{\partial^2 W}{\partial I_i\,\partial I_k} =
\dfrac{\partial^2 W}{\partial I_k\,\partial I_i}$ (теорема Клеро). Звідси негайно випливає теорема взаємності --- для \emphx{будь-якої} пари
контурів $i,k$ з довільної їх кількості:
\begin{equation*}
	L_{ik} = L_{ki}.
\end{equation*}

\begin{Attention}
	Ключова ідея доведення --- те, що енергія $W$ є \emphx{функцією стану} системи (потенціалом): вона не залежить від передісторії струмів, а
	лише від їхніх поточних значень. Саме тому до неї застосовна теорема про незалежність мішаних похідних від порядку диференціювання --- і саме
	вона й дає теорему взаємності. Це доведення одразу охоплює \emphx{будь-яку} кількість контурів і будь-яку пару індексів $i,k$ --- на відміну
	від прямого порівняння двох геометрично різних інтегралів (потоку контуру $i$ крізь контур $k$, і навпаки), яке саме собою нічого не підказує
	про їхню рівність. Існує й інше, суто геометричне доведення --- через явну формулу Неймана для взаємної індуктивності як подвійного
	контурного інтеграла, симетричного відносно перестановки контурів; воно спирається на векторний потенціал і виходить за межі цього курсу.
\end{Attention}

Практично корисний наслідок теореми: якщо по обох контурах течуть \emphx{однакові} струми $I$, то потоки, які вони наводять один в одному, теж
рівні, $\Phi_{12}=\Phi_{21}$, --- саме це використано в наступному прикладі.


\begin{Example}{Два коаксіальні кільця}
Два тонкі колові провідники зі спільною віссю лежать в одній площині; радіус зовнішнього --- $R_1$, внутрішнього --- $R_2$, причому $R_2\ll R_1$
(рис.~\ref{pic:CoaxialLoops}). Знайдемо магнітний потік крізь площу зовнішнього провідника, якщо внутрішнім тече струм $I$.

\begin{center}
	\localinput{coaxial-loops.tikz}
	\captionof{figure}{Два коаксіальні колові контури різних радіусів\label{pic:CoaxialLoops}}
\end{center}

Пряме обчислення потоку від маленького кільця крізь велике --- складна геометрична задача (поле маленького кільця складно неоднорідне на
масштабі великого). Натомість застосуємо \emphx{теорему взаємності}: обчислимо потік у \emphx{зворотному} напрямку --- від великого кільця крізь
маленьке, --- це набагато простіше, оскільки на масштабі маленького кільця поле великого можна вважати практично однорідним і рівним полю в
центрі великого кільця, $B_1 = 2\pi I/(cR_1)$ (гл.~\ref{MagFieldVacuum}). Тоді
\begin{equation*}
	\Phi_{12} = B_1\cdot \pi R_2^2 = \frac{2\pi^2 I R_2^2}{cR_1},
\end{equation*}
і, за теоремою взаємності, той самий потік (при тому самому струмі $I$) виникає й у прямій задачі:
\begin{equation*}
	\Phi_{21} = \Phi_{12} = \frac{2\pi^2 I R_2^2}{cR_1}.
\end{equation*}
\end{Example}

\begin{Example}{Обертовий магнітний диполь}
Магнітний диполь з моментом $p_m$ обертається з циклічною частотою $\omega$ навколо осі, що проходить через його центр перпендикулярно до $\vect
p_m$. На відстані $\ell \gg a$ від диполя розташоване нерухоме плоске колове кільце радіуса $a$ і опору $R$, нормаль якого $\vect n$
перпендикулярна до осі обертання диполя (рис.~\ref{pic:RotatingDipole}); самоіндукцією кільця знехтуємо. Знайдемо струм у кільці.

\begin{center}
	\localinput{rotating-dipole.tikz}
	\captionof{figure}{Магнітний диполь, що обертається, і нерухоме кільце на відстані $\ell$\label{pic:RotatingDipole}}
\end{center}

За формулою для поля далекого магнітного диполя~\eqref{eq:DipoleFieldB} (гл.~\ref{MagFieldVacuum}), у площині кільця нормальна складова поля
диполя пропорційна проєкції $p_m$ на напрям $\vect n$, тобто $\propto \sin(\omega t)$ (оскільки диполь обертається в площині, що містить
$\vect n$); потік крізь кільце виявляється синусоїдальним, $\Phi(t) \propto \sin(\omega t)$ з амплітудою, що спадає як $1/\ell^3$. ЕРС індукції
$\emf = -\frac1c d\Phi/dt \propto \omega\cos(\omega t)/\ell^3$, а шуканий струм $I(t) = \emf(t)/R$ --- гармонічно змінюється з тією
самою частотою $\omega$, з якою обертається сам диполь, і амплітудою, пропорційною $p_m\omega a^2/(cR\ell^3)$.
\end{Example}

\begin{Attention}
	Цей приклад --- не абстрактна вправа: саме так, у принципі, можна дистанційно \emphx{виявляти обертання} намагніченого тіла (наприклад,
	астрофізичного об'єкта з магнітним моментом) --- реєструючи індукований у нерухомій котушці змінний струм, частота якого дорівнює частоті
	обертання джерела поля.
\end{Attention}

%% --------------------------------------------------------
\section{Енергія системи зв'язаних контурів}
%% --------------------------------------------------------

Формула~\eqref{eq:MagEnergyFormula} стосується \emphx{одного} відокремленого контуру. Виведемо тепер явний
вираз для енергії системи довільної кількості $N$ магнітно зв'язаних контурів зі струмами $I_1,\dots,I_N$, спираючись на щойно доведену
теорему взаємності $L_{ik}=L_{ki}$ і на загальний вираз для потоку~\eqref{eq:FluxThroughIthLoop}.

Скористаємось при цьому загальним виразом для прирощення енергії~\eqref{eq:EnergyDifferential1}: він був виведений \emphx{ще до} доведення
теореми взаємності й ніяк не спирався на симетрію $L_{ik}$ (лише на закон збереження енергії), тому лишається справедливим і тепер, коли
симетрія вже встановлена, --- саме її бракувало раніше, щоб довести цей вираз до кінця. Розіб'ємо прирощення енергії~\eqref{eq:EnergyDifferential1}
на дві частини --- від самоіндукції кожного контуру окремо та від взаємної індукції кожної пари контурів:
\begin{equation}\label{eq:EnergyIncrementDecomposition}
	dW = dW_\text{с}+dW_\text{в}, \qquad dW_\text{с}=\frac1{c^2}\sum_{i=1}^{N} L_iI_i\,dI_i, \qquad
	dW_\text{в}=\frac1{c^2}\sum_{\substack{i,k=1\\i\neq k}}^{N} L_{ik}I_i\,dI_k.
\end{equation}
Доданок $dW_\text{с}$ інтегрується так само, як і для одного контуру у формулі~\eqref{eq:MagEnergyFormula}:
\begin{equation}\label{eq:SelfInductionEnergy}
	W_\text{с}=\frac1{2c^2}\sum_{i=1}^{N} L_iI_i^2.
\end{equation}
У сумі $dW_\text{в}$ згрупуємо доданки попарно: для кожної пари індексів $i<k$ у подвійній сумі за $i\neq k$ трапляються рівно два доданки,
$L_{ik}I_i\,dI_k$ і $L_{ki}I_k\,dI_i$:
\begin{equation*}
	dW_\text{в} = \frac1{c^2}\sum_{i<k}\left(L_{ik}I_i\,dI_k+L_{ki}I_k\,dI_i\right).
\end{equation*}
Саме тут, і лише тут, потрібна теорема взаємності: підставивши $L_{ki}=L_{ik}$, кожна дужка перетворюється на повний диференціал добутку
$I_iI_k$,
\begin{equation*}
	L_{ik}I_i\,dI_k+L_{ik}I_k\,dI_i = L_{ik}\,d(I_iI_k),
\end{equation*}
звідки
\begin{equation*}
	dW_\text{в} = \frac1{c^2}\sum_{i<k}L_{ik}\,d(I_iI_k) = \frac1{2c^2}\sum_{\substack{i,k=1\\i\neq k}}^{N}L_{ik}\,d(I_iI_k),
\end{equation*}
і, інтегруючи від нульових струмів,
\begin{equation}\label{eq:MutualEnergyFinal}
	W_\text{в}=\frac1{2c^2}\sum_{\substack{i,k=1\\i\neq k}}^{N}L_{ik}I_iI_k.
\end{equation}
Об'єднуючи~\eqref{eq:SelfInductionEnergy} і~\eqref{eq:MutualEnergyFinal}, отримуємо повну енергію системи контурів:
\begin{equation}\label{eq:MultiLoopEnergy}
	\tcbhighmath{
		W = W_\text{с}+W_\text{в} = \frac1{2c^2}\sum_{i,k=1}^{N} L_{ik}I_iI_k = \frac1{2c}\sum_{i=1}^{N} I_i\Phi_i,
	}
\end{equation}
де в останній рівності знову використано означення потоку~\eqref{eq:FluxThroughIthLoop}. Саме завдяки симетрії $L_{ik}=L_{ki}$ подвійну
суму~\eqref{eq:MultiLoopEnergy} можна писати компактно, не розрізняючи порядок індексів у кожному доданку, --- без теореми взаємності довелося б
щоразу стежити, який із двох, узагалі кажучи різних, коефіцієнтів стоїть у кожному доданку.

Для системи з двох контурів формула~\eqref{eq:MultiLoopEnergy} набуває вигляду
\begin{equation}\label{eq:TwoLoopEnergy}
	W = \frac1{2c^2}\left(L_1I_1^2+L_2I_2^2+2L_{12}I_1I_2\right),
\end{equation}
де $L_1\equiv L_{11}$, $L_2\equiv L_{22}$ --- власні індуктивності контурів.

%% --------------------------------------------------------
\subsection*{Обмеження на взаємну індуктивність}
%% --------------------------------------------------------

Вираз~\eqref{eq:TwoLoopEnergy} для енергії двох контурів --- квадратична форма струмів $I_1,I_2$. Енергія магнітного поля не може бути
від'ємною за жодних значень струмів, $W\ge0$ (стан зі струмом не може мати меншу енергію, ніж стан без струму, --- інакше поле саме собою
\enquote{постачало} б енергію), тобто ця квадратична форма повинна бути додатно визначеною. Записавши~\eqref{eq:TwoLoopEnergy} у стандартному
вигляді $W=aI_1^2+2bI_1I_2+dI_2^2$ з $a=L_1/2c^2$, $b=L_{12}/2c^2$, $d=L_2/2c^2$, отримуємо звичайну умову додатної визначеності $b^2\le ad$,
тобто
\begin{equation}\label{eq:InductanceInequality}
	\tcbhighmath{
		L_{12}^2 \le L_1L_2.
	}
\end{equation}
Фізично ця нерівність означає, що магнітний зв'язок між двома контурами не може бути \enquote{сильнішим}, ніж дозволяють їхні власні
індуктивності: частина енергії, яку створює струм одного контуру, завжди залишається в його \emphx{власному} магнітному полі, і лише частина
--- у полі, зчепленому з іншим контуром. Рівність $L_{12}^2=L_1L_2$ означає \emphx{ідеальний} магнітний зв'язок (весь потік одного контуру
повністю зчеплений з іншим, без розсіювання) --- граничний, майже не досяжний на практиці випадок, до якого, однак, впритул наближаються добре
сконструйовані обмотки трансформатора на замкненому феромагнітному осерді (наступний параграф).

%% --------------------------------------------------------
\subsection*{Вираз енергії системи контурів через поля}
%% --------------------------------------------------------

	Той самий результат~\eqref{eq:TwoLoopEnergy} можна отримати зовсім іншим шляхом --- не через роботу джерел, а через \emphx{локалізацію
	енергії в полі}, узагальнюючи формулу~\eqref{eq:MagEnergyDensityGeneral} на випадок поля, породженого одразу кількома контурами (у лінійному
	середовищі з $\mu=\const$, як і в прикладі з двома обмотками вище). Нехай $\Bfield_1$ --- поле струму $I_1$, а $\Bfield_2$ --- поле струму
	$I_2$; за принципом суперпозиції $\Bfield=\Bfield_1+\Bfield_2$, і
	\begin{equation*}
		W=\iiint\limits_V \frac{B^2}{8\pi\mu}\,dV = \iiint\limits_V \frac{B_1^2}{8\pi\mu}\,dV+\iiint\limits_V \frac{B_2^2}{8\pi\mu}\,dV+
		\iiint\limits_V \frac{\Bfield_1\cdot\Bfield_2}{4\pi\mu}\,dV.
	\end{equation*}
	Перші два інтеграли --- це \enquote{власні} енергії кожного контуру окремо ($\tfrac12L_1I_1^2/c^2$ і $\tfrac12L_2I_2^2/c^2$), а третій,
	\enquote{перехресний}, доданок і дає взаємну енергію $L_{12}I_1I_2/c^2$ з формули~\eqref{eq:TwoLoopEnergy}; для довільної кількості контурів
	міркування те саме, лише з подвійною сумою перехресних доданків замість одного.

	Це трактування відразу дає кілька наслідків, не таких очевидних із формули роботи джерел:
	\begin{enumerate}
		\item Енергія $W$ --- \emphx{завжди додатна} за будь-яких струмів, оскільки під інтегралом стоїть невід'ємна величина $B^2$.
		\item Енергія системи контурів \emphx{не адитивна}: через перехресний доданок вона не дорівнює сумі енергій окремих контурів.
		\item Коефіцієнт при $I_1I_2$ у третьому інтегралі --- тобто сам інтеграл $\iiint \Bfield_1\cdot\Bfield_2/4\pi\mu\,dV$ --- очевидно
		      \emphx{симетричний} щодо перестановки індексів 1 і 2 (скалярний добуток не залежить від порядку множників). Тому рівність
		      $L_{12}=L_{21}$ тут випливає не з абстрактного доведення через мішані похідні енергії (яке ми навели раніше), а як \emphx{пряме
		      геометричне} твердження про симетрію інтеграла --- друге, незалежне доведення теореми взаємності.
		\item Звідси ж --- і друге означення взаємної індуктивності, через розподіл полів у просторі, а не через потік крізь контур:
		      \begin{equation}\label{eq:MutualInductanceViaFieldIntegral}
			      L_{12}=\frac{c^2}{I_1I_2}\iiint\limits_V \frac{\Bfield_1\cdot\Bfield_2}{4\pi\mu}\,dV.
		      \end{equation}
	\end{enumerate}
	Отже, обидва підходи рівноправні й дають той самий результат, що ще раз
	підтверджує: магнітна енергія --- це справді енергія \emphx{поля}, розподілена в просторі, а не якась умовна властивість самих контурів.

\begin{Attention}
	Та сама картина --- додатність повного інтеграла й неадитивність через перехресний доданок --- вже траплялась нам для електричного поля
	(гл.~\ref{ElEnergy}): розклад $E^2=E_1^2+E_2^2+2\Efield_1\cdot\Efield_2$ так само давав завжди додатну повну енергію
	\eqref{eq:GeneralEnergyDensity} проти знакозмінної \enquote{дискретної} взаємної енергії~\eqref{eq:SelfMutualEnergy}. Різниця лише в тому,
	що там власна енергія точкового заряду розбігається, тому в \enquote{дискретній} формулі~\eqref{eq:PairEnergy} лишається саме
	взаємний доданок; тут же власні енергії контурів $L_iI_i^2/2c^2$ скінченні й додатні самі собою, а знак може мати лише перехресний доданок
	$L_{ik}I_iI_k/c^2$ --- і все ж повна енергія додатна завжди, бо додатна сама величина $B^2$ під інтегралом.

\end{Attention}

%% --------------------------------------------------------
\section{Принцип роботи трансформаторів}
%% --------------------------------------------------------

Явище взаємної індукції лежить в основі роботи \emphx{трансформаторів} --- пристроїв для підвищення чи зниження напруги змінного струму
(рис.~\ref{pic:Transformer}). Первинна й вторинна обмотки з $N_1$ і $N_2$ витками укріплені на спільному замкненому феромагнітному осерді;
завдяки великому $\mu$ феромагнетика (гл.~\ref{MagFieldMedia}) магнітний потік майже повністю локалізований в осерді й пронизує обидві обмотки
однаковим потоком $\Phi$. Це --- практична реалізація того самого \emphx{ідеального} магнітного зв'язку, граничного випадку $L_{12}^2=L_1L_2$ з
попереднього параграфа: розсіюванням потоку поза осердям тут свідомо нехтують, а отже, увесь потік, породжений первинною обмоткою, зчеплений і
з вторинною. Змінна ЕРС $\emf_1$ у первинній обмотці створює струм $I_1$ і змінний магнітний потік $\Phi$; цей потік пронизує й витки вторинної
обмотки, наводячи в ній ЕРС взаємної індукції, а в первинній --- ЕРС самоіндукції.

\begin{figure}[h!]\centering
	\includegraphics[width=0.5\linewidth]{Transformator}
	\caption{Принципова схема трансформатора\label{pic:Transformer}}
\end{figure}

Закон Ома для первинного кола: $\emf_1 - \frac1c\dfrac{d}{dt}(N_1\Phi) = I_1R_1$. У режимі швидкозмінних полів (типовому для трансформаторів)
$I_1R_1 \ll \emf_1,\ \tfrac{d}{dt}(N_1\Phi)$, тому наближено
\begin{equation}\label{eq:TransformerEMF1}
	\emf_1 \approx \frac{N_1}c\frac{d\Phi}{dt},
\end{equation}
а ЕРС у вторинній обмотці, оскільки потік $\Phi$ в осерді той самий,
\begin{equation}\label{eq:TransformerEMF2}
	\emf_2 = -\frac{N_2}c\frac{d\Phi}{dt} = -\frac{N_2}{N_1}\emf_1
\end{equation}
(знак \enquote{мінус} означає протилежність фаз). Відношення
\begin{equation}\label{eq:TransformationRatio}
	\tcbhighmath{
		\eta = \frac{N_2}{N_1}
	}
\end{equation}
називають \emphx{коефіцієнтом трансформації}: якщо $\eta>1$ --- трансформатор \emphx{підвищувальний} (збільшує ЕРС), якщо $\eta<1$ ---
\emphx{знижувальний} (зменшує ЕРС).

А як перетворюється струм? Тут якраз і працює закон збереження енергії. У наближенні ідеального трансформатора --- без омічних втрат
($R_1,R_2\to0$) і без розсіювання потоку --- в обмотках не накопичується енергія понад ту, що в кожен момент проходить крізь них: уся
потужність, яку джерело вкладає в первинне коло, миттєво передається через спільне поле в коло навантаження вторинної обмотки, тобто
потужність зберігається,
\begin{equation*}
	I_1\emf_1 = -I_2\emf_2
\end{equation*}
(знак \enquote{мінус} --- та сама протифазність, що й у формулі~\eqref{eq:TransformerEMF2}: коли первинна обмотка \emphx{споживає} потужність
від джерела, вторинна \emphx{віддає} її навантаженню). Підставляючи~\eqref{eq:TransformerEMF2}, знаходимо коефіцієнт трансформації \emphx{струмів}:
\begin{equation}\label{eq:TransformerCurrentRatio}
	\tcbhighmath{
		\frac{I_2}{I_1} = -\frac{\emf_1}{\emf_2} = \frac1\eta.
	}
\end{equation}
Отже, у скільки разів трансформатор підвищує напругу, у стільки ж разів він знижує струм, --- і навпаки: закон збереження енергії не дозволяє
одночасно виграти і в напрузі, і в струмі. Саме тому передавати електроенергію на великі відстані вигідно при високій напрузі й, відповідно,
малому струмі (а отже, і з малими тепловими втратами $I^2R$ у проводах лінії передачі), а споживачам подавати вже при низькій, безпечній
напрузі --- і саме взаємна індукція, а не який-небудь прямий електричний контакт, дозволяє так легко узгоджувати ці дві вимоги.

%% --------------------------------------------------------
\section{Пондеромоторні сили в магнітному полі}
%% --------------------------------------------------------


%% --------------------------------------------------------
\subsection*{Метод віртуальних переміщень для магнітних сил}
%% --------------------------------------------------------

Так само, як і в електростатиці (гл.~\ref{ElEnergy}), силу, що діє на провідник чи магнетик з боку магнітного поля, зручно шукати методом
віртуальних переміщень, спираючись на закон збереження енергії. Робота джерел струму $\delta A_\text{дж}$ витрачається на теплоту $\delta Q$, на
приріст магнітної енергії $dW$ і на механічну роботу $\delta A_\text{мех}$:
\begin{equation}\label{eq:PonderomotiveWork}
	\delta A_\text{дж} = \delta Q + dW + \delta A_\text{мех}.
\end{equation}
З іншого боку, за законом Ома, робота джерела дорівнює $\delta A_\text{дж} = \delta Q - \emf_\text{сі}I\, dt$ (ЕРС самоіндукції --- частина
роботи джерела, витрачена на подолання цієї ЕРС). Порівнюючи два вирази для $\delta A_\text{дж}$:
\begin{equation*}
	-\emf_\text{сі}I\, dt = dW + \delta A_\text{мех}.
\end{equation*}

\emphx{Випадок $I=\const$} (струм у колі підтримується сталим джерелом). Тоді, за формулою~\eqref{eq:SelfInductionEMF}, $-\mathcal
E_\text{сі}I\,dt = \dfrac{I^2}{c^2}dL$, тоді як приріст магнітної енергії за формулою~\eqref{eq:MagEnergyFormula} становить лише половину цього,
$dW = \dfrac{I^2}{2c^2}dL$ (оскільки $I=\const$, а не $\Phi$). Отже, $\delta A_\text{мех} = +dW$, і сила
\begin{equation}\label{eq:ForceConstI}
	\tcbhighmath{
		F_x = +\left(\frac{dW}{dx}\right)_I.
	}
\end{equation}

\emphx{Випадок $\Phi=\const$} (наприклад, контур з надпровідника, де потік зберігається за формулою~\eqref{eq:FluxConservationZeroR}). Тоді
$\emf_\text{сі} = 0$ (адже $d\Phi/dt=0$), тому $\delta A_\text{мех} = -dW$, і
\begin{equation}\label{eq:ForceConstPhi}
	\tcbhighmath{
		F_x = -\left(\frac{dW}{dx}\right)_\Phi.
	}
\end{equation}

\begin{Attention}
	Знаки формул~\eqref{eq:ForceConstI} і~\eqref{eq:ForceConstPhi} --- протилежні до відповідних формул для електричної
	сили,~\eqref{eq:ForceConstQ} і~\eqref{eq:ForceConstPhi} з гл.~\ref{ElEnergy} (там при сталому заряді був \enquote{мінус}, тут при сталому
	струмі --- \enquote{плюс}, і навпаки). Причина --- у формулах~\eqref{eq:SelfInductionFlux} і~\eqref{eq:MagEnergyFormula}: магнітна енергія
	пропорційна $LI^2$, а не $q^2/C$, тому при диференціюванні за $I=\const$ (аналог $\Delta\phi=\const$ в електриці) чи $\Phi=\const$ (аналог
	$q=\const$) знаки \enquote{переставляються} місцями порівняно з електричним випадком --- варто бути особливо уважним і щоразу передивлятись
	це виведення заново, а не покладатись на \enquote{інтуїтивну} аналогію зі знаками.
\end{Attention}

%% --------------------------------------------------------
\subsection*{Загальний енергетичний метод: вільна енергія}
%% --------------------------------------------------------

Розглянутий метод неявно передбачав, що при віртуальному переміщенні жодні інші, немагнітні властивості речовини не змінюються. Строго кажучи,
закон збереження енергії, на якому ґрунтується метод віртуальних переміщень, застосовний за \emphx{сталої температури} $T=\const$ (щоб не
враховувати змін властивостей речовини, спричинених неелектромагнітними чинниками); а тоді, згідно з термодинамічним співвідношенням
$dF=-S\,dT-P\,dV$, роботу сил поля визначає не повна, а \emphx{вільна енергія} $F$. Для лінійного середовища (коли індуктивності не залежать від
струмів) вільна енергія системи контурів чисельно збігається з уже знайденою магнітною енергією~\eqref{eq:MultiLoopEnergy}:
\begin{equation}\label{eq:FreeEnergyDefinition}
	F = \frac1{2c}\sum_{i=1}^{N} I_i\Phi_i \qquad(=W).
\end{equation}
Переконаймося узагальненим методом віртуальних переміщень, застосованим одразу до системи $N$ контурів, що результат дійсно узгоджується з
формулами~\eqref{eq:ForceConstI} і~\eqref{eq:ForceConstPhi}, отриманими вище для одного контуру.

\paragraph{1) Сталі потоки.} Якщо опори контурів знехтувано малі, магнітні потоки при віртуальному переміщенні не змінюються. Справді, за законом Ома з урахуванням ЕРС
індукції, $I_iR_i = -\dfrac1c\dfrac{d\Phi_i}{dt}$; якщо $R_i\to0$ при скінченному струмі $I_i$, ліва частина прямує до нуля, а отже
$d\Phi_i/dt=0$ --- тобто потік $\Phi_i = \const$ лишається сталим у кожному з $N$ контурів системи. ЕРС індукції тоді не виникає, і вся робота сил поля
виконується за рахунок спаду вільної енергії:
\begin{equation}\label{eq:WorkConstFlux}
	\delta A_\text{поле} = f\,\delta x = -\left(dF\right)_{\{\Phi_i\},T}, \qquad \left(dF\right)_{\{\Phi_i\},T} = \frac1{2c}\sum_{i=1}^{N}\Phi_i\,dI_i,
\end{equation}
де останній вираз --- диференціал $F=\frac1{2c}\sum_i I_i\Phi_i$, узятий при сталих $\Phi_i$.

\paragraph{2) Сталі струми.} Нехай тепер у кожен контур включено батарею ЕРС $\emf_i^\text{бат}$, що врівноважує ЕРС індукції, яка виникає
при будь-якій зміні потоків: $\emf_i^\text{бат} = -\emf_i^\text{інд} = \dfrac1c\dfrac{d\Phi_i}{dt}$. За час віртуального переміщення
$dt$ батареї виконують роботу
\begin{equation*}
	\delta A_\text{бат} = \sum_{i=1}^{N} I_i\emf_i^\text{бат}\,dt = \frac1c\sum_{i=1}^{N} I_i\,d\Phi_i,
\end{equation*}
яка йде і на роботу сил поля, і на зміну вільної енергії: $\delta A_\text{бат} = \delta A_\text{поле} + \left(dF\right)_{\{I_i\},T}$. Диференціал
$F$ при сталих струмах, $\left(dF\right)_{\{I_i\},T} = \frac1{2c}\sum_i I_i\,d\Phi_i$, становить рівно половину роботи батарей, тому
\begin{equation}\label{eq:WorkConstCurrent}
	\delta A_\text{поле} = \delta A_\text{бат} - \left(dF\right)_{\{I_i\},T} = \frac1c\sum_{i}I_i\,d\Phi_i-\frac1{2c}\sum_iI_i\,d\Phi_i =
	\frac1{2c}\sum_i I_i\,d\Phi_i = +\left(dF\right)_{\{I_i\},T}.
\end{equation}

Порівнюючи~\eqref{eq:WorkConstFlux} і~\eqref{eq:WorkConstCurrent} з формулами~\eqref{eq:ForceConstI} і~\eqref{eq:ForceConstPhi}, бачимо, що
останні --- окремий випадок ($N=1$, $F=W$) загального результату
\begin{equation}\label{eq:WorkBothApproaches}
	\tcbhighmath{
		\delta A_\text{поле} = f\,\delta x = -\left(dF\right)_{\{\Phi_i\},T} = +\left(dF\right)_{\{I_i\},T}.
	}
\end{equation}

\begin{Attention}
	Розрізняти $W$ і $F$ важливо лише тоді, коли властивості середовища (наприклад, магнітна проникність $\mu$) залежать від температури, а
	віртуальне переміщення супроводжується тепловими ефектами. У прикладах цього розділу $\mu=\const$, тому $F$ і $W$ --- буквально та сама
	величина, і формула~\eqref{eq:WorkBothApproaches} нічого не змінює по суті порівняно з~\eqref{eq:ForceConstI},~\eqref{eq:ForceConstPhi} ---
	вона лише узагальнює їх одразу на систему довільної кількості контурів і явно фіксує умову $T=\const$, яка неявно малась на увазі й раніше.
\end{Attention}


\begin{Example}{Втягування стержня з магнетика в соленоїд}

Знайдемо, користуючись загальним методом, силу, з якою втягується у вакуумний соленоїд довгий циліндричний стержень з магнітною проникністю
$\mu$ і площею поперечного перерізу $S$. По соленоїду, що має густину намотування $n=N/\ell$, циркулює струм $I$.


\begin{center}%[h!]\centering
	\includegraphics[width=0.55\linewidth]{pull_in_solenoid.jpg}
	\captionof*{figure}{Осердя втягується в котушку зі струмом --- приклад пондеромоторної сили в
		магнітному полі\label{pic:RodInSolenoid}}
\end{center}


%\begin{center}%[h!]\centering
%%	\localinput{rod-in-solenoid.tikz}
%	\captionof*{figure}{Стержень з магнетика, введений у соленоїд, зміщується праворуч на малу відстань $\delta x$\label{pic:RodInSolenoid}}
%\end{center}

Вважатимемо, що струм у соленоїді підтримується сталим, тому користуємось формулою при сталому струмі,
\begin{equation}\label{eq:WorkConstCurrentRod}
	\delta A_\text{поле} = f\,\delta x = \left(dF\right)_{I,T} = \frac{I}{2c}\,d\Phi
\end{equation}
(тут один контур, $N=1$, тому $F=W=I\Phi/2c$, як і в означенні~\eqref{eq:MagEnergyFormula}). Далі, оскільки соленоїд довгий, а точки
спостереження обрано далеко від межі розділу стержень--вакуум і від кінців соленоїда, напруженість поля там усюди та сама, що й в ідеальному
нескінченному соленоїді: $H=4\pi ni/c=4\pi nI/c$ (гл.~\ref{MagFieldMedia}), звідки $I=cH/4\pi n$, а отже
\begin{equation}\label{eq:FreeEnergyDiffViaH}
	\left(dF\right)_{I,T} = \frac1{8\pi n}H\,d\Phi.
\end{equation}

Знайдемо зміну потоку $d\Phi$ при зміщенні стержня на $\delta x$. Потік $\Phi=BS$ (через один виток) змінюється за рахунок того, що частина
довжини соленоїда, яка досі містила вакуум ($\mu=1$), заповнюється речовиною з проникністю $\mu$: при зміщенні стержня на $\delta x$ число
витків, що охоплюють вакуумну частину, зменшується на $n\,\delta x$, а число витків, що охоплюють частину з магнетиком, на стільки ж
збільшується. Тому, нехтуючи неоднорідністю поля поблизу межі розділу й кінців соленоїда (об'єм цих ділянок при малому $\delta x$ практично не
змінюється) і використовуючи $B_\text{стержень}=\mu H$, $B_\text{вакуум}=H$ (гл.~\ref{MagFieldMedia}):
\begin{equation}\label{eq:FluxChangeRod}
	d\Phi = n\,\delta x\left(B_\text{стержень}-B_\text{вакуум}\right)S = n\,\delta x\,(\mu-1)HS.
\end{equation}
Підставляючи~\eqref{eq:FluxChangeRod} у~\eqref{eq:FreeEnergyDiffViaH}, знаходимо
\begin{equation*}
	\left(dF\right)_{I,T} = \frac{\mu-1}{8\pi}H^2S\,\delta x,
\end{equation*}
звідки, за формулою~\eqref{eq:WorkConstCurrentRod},
\begin{equation}\label{eq:ForceOnRod}
		f = \frac{\left(dF\right)_{I,T}}{\delta x} = \frac{\mu-1}{8\pi}H^2S.
\end{equation}
При $\mu>1$ (парамагнетик) сила додатна --- стержень втягується в соленоїд; при $\mu<1$ (діамагнетик) сила від'ємна --- стержень виштовхується
з нього.
\end{Example}

\begin{Attention}
	Формула~\eqref{eq:ForceOnRod} --- ще одне підтвердження загального правила: парамагнетики ($\mu>1$) втягуються в область сильнішого поля, а
	діамагнетики ($\mu<1$) --- виштовхуються з неї. Сила пропорційна $H^2$ (тобто не залежить від напрямку поля) і площі перерізу стержня $S$ ---
	чим \enquote{товщий} стержень, тим більша втягувальна сила. Порівняйте цей результат із формулою~\eqref{eq:ElectromagnetForce} для підйомної
	сили електромагніта: там сила теж виявилась пропорційною $B^2S$ (з точністю до множника, зумовленого іншою геометрією задачі), що ще раз
	наголошує на універсальності енергетичного методу --- він дає змогу знаходити силу в найрізноманітніших конфігураціях, не вдаючись щоразу до
	прямого обчислення сил Ампера.
\end{Attention}

\begin{Example}{Підйомна сила електромагніта}
Розглянемо електромагніт --- П-по\-діб\-не осердя з феромагнетика проникності $\mu$, обмотане $N$ витками зі струмом $I$, до полюсів якого прикладено
якір (пластину), відокремлений повітряним зазором завтовшки $x$ з кожного боку. Знайдемо силу притягання якоря.

\begin{center}%[1.5][h!]\centering
	\localinput{electromagnet.tikz}
	\captionof*{figure}{Електромагніт з якорем, відокремленим повітряним зазором $x$\label{pic:ElectromagnetForce}}
\end{center}

За теоремою про циркуляцію для $\Hfield$, застосованою до контуру, що проходить через осердя (довжиною $\ell$) і двічі перетинає зазор (загальною
товщиною $2x$):
\begin{equation*}
	\frac{B}{\mu}\ell + B\cdot 2x = \frac{4\pi}{c}IN \quad\Rightarrow\quad B = \frac{4\pi\mu}{c}\,\frac{NI}{\ell+2\mu x}
\end{equation*}
(у повітряному зазорі $\mu_\text{пов}\approx1$, тому там $\Hfield=\Bfield$; поле в зазорі й осерді те саме завдяки неперервності $B_n$ на межі,
гл.~\ref{MagFieldMedia}). Магнітний потік крізь переріз $S$ магніту: $\Phi = NBS = \dfrac{4\pi\mu}{c}\dfrac{N^2I}{\ell+2\mu x}S$.

За умови сталого струму ($I=\const$) зміна енергії поля при віртуальній зміні зазору $dW = \dfrac{I\, d\Phi}{2c}$ (з формули~\eqref{eq:MagEnergyFormula}, $W=I\Phi/(2c)$ при $I=\const$), і, за формулою~\eqref{eq:ForceConstI},
\begin{equation*}
	F_x = \left(\frac{dW}{dx}\right)_I = \frac{I}{2c}\left(\frac{d\Phi}{dx}\right)_I = -\frac{4\pi\mu^2}{c^2}\,\frac{N^2I^2}{(\ell+2\mu x)^2}S =
	-\frac{B^2}{8\pi}\cdot 2S.
\end{equation*}
Отже,
\begin{equation}\label{eq:ElectromagnetForce}
		F_x = -\frac{B^2}{4\pi}S.
\end{equation}
Від'ємний знак означає, що сила --- сила \emphx{притягання}: якір втягується в зазор, зменшуючи $x$, --- цілком у дусі загального висновку про
те, що магнетики (тут --- феромагнітний якір) завжди втягуються в область сильнішого поля (гл.~\ref{MagFieldVacuum}).
\end{Example}


%% --------------------------------------------------------
\section*{Підсумки}
%% --------------------------------------------------------


Шлях цього розділу був обернений до того, яким вводилась енергія електричного поля: там вона народжувалась \emphx{напряму} з роботи проти сили Кулона. Тут так не вийшло: магнітна сила, за Лоренцем, роботи над зарядами не виконує ніколи --- ні для сталого поля, ні для змінного. Але, змінюючись, магнітне поле породжує інше поле --- \emphx{вихрове електричне} --- і вже воно виконує роботу над зарядами. Саме цю роботу, на подолання ЕРС \emphx{самоіндукції}, і виконує джерело струму, а не роботу проти сили Ампера.

З цієї роботи ми й отримали густину енергії $w=BH/8\pi$ --- за формою тотожну електростатичній $w=ED/8\pi$ (гл.~\ref{ElEnergy}). Та сама логіка лежить і в основі пондеромоторних сил (втягування магнетика в соленоїд, притягання якоря електромагніта). На макрорівні здається, що механічну роботу виконує саме магнітна сила. Проте для заряду, що рухається \emphx{разом} із провідником, ця ж сила одночасно гальмує дрейфовий струм. Робота магнітної сили над макроскопічним рухом тіла в точності компенсується її ж від'ємною роботою над струмом --- так що сумарно, як і належить, сила Лоренца роботи не виконує. У мові ж електродинамічних контурів рух тіла змінює індуктивність $L(x)$, викликаючи додаткову ЕРС індукції. Саме джерело струму, долаючи цю ЕРС, компенсує гальмування носіїв і здійснює підсумкову роботу, яка ділиться порівну: одна її частина йде на поповнення енергії поля з густиною $w = BH/8\pi$, а інша --- переходить у механічний розгін тіла.

Звідси головний висновок розділу: магнітне поле --- не \emphx{виконавець}, а \emphx{посередник}. Енергію воно \emphx{зберігає}, але обмінюється нею із зарядами лише через електричне поле, яке саме ж і породжує, змінюючись в часі.


%З цієї роботи ми й отримали густину енергії $w=BH/8\pi$ --- за формою тотожну електростатичній $w=ED/8\pi$ (гл.~\ref{ElEnergy}). Тієї самої природи --- посередницької --- виявились і пондеромоторні сили: вони теж не породжуються магнітним полем безпосередньо, а знаходяться через ту саму енергію поля, запасену й вивільнювану через електричне.





\begin{Questions}
\begin{quizlist}
	\item Виведіть енергію магнітного поля контуру зі струмом $W = LI^2/2c^2$ з роботи джерела ЕРС. На що йде решта роботи джерела?
	\item Як звідси отримати густину енергії магнітного поля $w = BH/8\pi$ для соленоїда? Як це виведення узагальнюється на довільне поле і яку роль у ньому відіграє вихрове електричне поле?
	\item Дайте означення коефіцієнтів взаємної індуктивності $L_{12}$, $L_{21}$ (чим вони відрізняються від власної індуктивності) і покажіть на прикладі двох обмоток на спільному осерді, що $L_{12}=L_{21}$. За яких умов ця рівність може порушуватися і яку роль тут відіграє феромагнітне осердя?
	\item Сформулюйте теорему взаємності. Наведіть її доведення через частинні похідні енергії за струмами і теорему Клеро. Чому це доведення одразу узагальнюється на будь-яку кількість контурів? Як теорема спрощує обчислення потоку, як у прикладі з двома коаксіальними кільцями?
	\item Узагальніть формулу енергії одного контуру на систему $N$ зв'язаних контурів. Чому саме теорема взаємності дозволяє писати подвійну суму $W=\frac1{2c^2}\sum L_{ik}I_iI_k$, не турбуючись про порядок індексів? Звідки береться нерівність $L_{12}^2\le L_1L_2$ і що вона означає фізично?
	\item Поясніть принцип дії трансформатора. Що таке коефіцієнт трансформації, чим відрізняються підвищувальний і знижувальний трансформатори і чому передавати електроенергію вигідно при високій напрузі?
	\item Виведіть формули для пондеромоторної сили при $I=\const$ та при $\Phi=\const$. Чому їхні знаки протилежні до відповідних формул електростатики?
	\item Що таке вільна енергія в контексті енергетичного методу обчислення магнітних сил і чому саме вона, а не повна енергія, визначає роботу поля? Знайдіть силу, з якою парамагнітний чи діамагнітний стержень втягується (виштовхується) із соленоїда.
	\item Знайдіть підйомну силу електромагніта (рис.~\ref{pic:ElectromagnetForce}). Чому вона є силою притягання?
\end{quizlist}
\end{Questions}